% !TeX root = ../../Book3.tex
\chapter*{序言}
\addcontentsline{toc}{chapter}{序言}
\markboth{序言}{序言}

设 \(k\) 为一个域。在 \(k\) 中，方程 \(x=0\) 与 \(x^2=0\) 有相同的解；若只记下零点集，它们之间似乎没有区别。但相应的商环 \(k[x]/(x)\) 与 \(k[x]/(x^2)\) 不同：在后一个环中，\(x\) 的剩余类不是零，它的平方却是零。几何上的一个点，原来还可以带着零点集看不见的信息。第三卷从这样的现象出发，研究多项式方程、交换环和几何对象之间究竟能建立怎样的联系。

\ProseParagraphBreak
这些问题有很长的来历。十九世纪以来，对多项式方程的研究逐渐不满足于求出几个具体的根。Hilbert 基定理在适当的有限性假设下保证多项式理想有限生成，零点定理在代数闭域上把理想与公共零点联系起来；Noether 性处理有限性，Zariski 拓扑则把零点组织成空间。一个方程可以画出曲线，许多方程可以确定交集；把方程组成理想以后，我们还能问：它是否有几个分支，某处是否重合得更紧，以及函数在一个点附近是否仍可定义。

\ProseParagraphBreak
从整体走向局部，是本卷经常采用的办法。研究一条曲线在某点附近的性质时，我们允许那些在该点不为零的函数充当分母，得到局部环。不同的点会给出不同的局部环；把它们放回一起，又要问局部所得的信息能否拼成整体结论。这种往返并不总是顺利。一个模在每个点附近自由，不一定能选出一组全局的基；一个空间的点都已数清，也未必就看清了它上面的函数。

\ProseParagraphBreak
第一编从素理想、素谱和局部化开始，继而研究有限性、准素分解与整扩张。Noether 正规化引理和 Hilbert 零点定理使多项式环中的代数问题进入仿射几何。读到坐标环时，可以回想开头的两个方程：只看点会丢失什么，保留环又多知道了什么？尖点、结点和加厚点给出比口头解释更明确的答案。

\ProseParagraphBreak
第二编把存在性结论与计算接起来。Gröbner 基和消去法使一些理想运算化为具体多项式的计算；素理想链与 Hilbert 函数从不同方面测量维数。随后讨论离散赋值环、Dedekind 整环、平坦性、完备化和 Kähler 微分。它们看上去各有用途，却常在同一个问题中相遇：先计算方程的关系，再考察一点附近的阶数、切空间，以及参数变化时哪些性质仍能保持。

\ProseParagraphBreak
第三编进一步研究局部环。维数相同的两个局部环，性质可能相差很远；正则序列、深度和自由消解帮助我们说明差别发生在哪里。Koszul 复形、\(\operatorname{Ext}\) 与 \(\operatorname{Tor}\) 在这里都服务于具体的环与模。正则局部环、Cohen--Macaulay 性和正规性分别从不同角度约束局部结构，分次消解则把同调信息带回可计算的多项式关系。

\ProseParagraphBreak
第四编的三章均供选读。局部上同调和对偶把局部结构继续向前推进；Hensel 提升、光滑与平展讨论方程的解及其随基底变化的行为；最后从层和仿射概形出发，走到 Proj 与射影直线。概形部分只介绍基本构造和少量例子。一般定理若需借助本书尚未证明的结果，正文会说明所用的条件与来源。

\ProseParagraphBreak
本卷适合已经熟悉理想、商环、域扩张，以及模、正合列和张量积的读者。第一卷的环论与第二卷前七章提供主要基础；不必先读完第一卷的高级 Galois 理论或第二卷的非交换专题。前十四章可以先行学习。进入第十五至二十一章时，若尚不熟悉复形、消解、\(\operatorname{Ext}\) 和 \(\operatorname{Tor}\)，可先补第四卷第五至八章（含第八章第五节）的基本内容，再回来读深度与局部结构；不必等到学完导出范畴。第四编各章所需基础并不完全相同，可按编导读选择。

\ProseParagraphBreak
附录也不必依字母顺序读。想看整性与赋值怎样进入算术，可转向数域整数环；关心计算，可在相应正文之后阅读局部标准基与实根问题；微分代数则从导子和微分多项式另开一条线。每篇附录开头都说明所需基础和讨论范围，不必为了读正文先把附录全部读完。

\ProseParagraphBreak
学习交换代数，手边留几个反复使用的例子很有帮助：多项式环、一个含幂零元的商环、一条奇异曲线的坐标环，以及一点处的局部环。读到新的性质，可以先在这些对象上检验它；读到一个反例，也应追问究竟是哪条假设失效。计算 Gröbner 基时保存约化等式，使用局部化时写清分母，比较两个模时亲自构造映射。定义与定理往往只有几行，理解它们却离不开这些具体核对。

\ProseParagraphBreak
各章给出学习目标、习题与解答，并在需要处标明进一步阅读的来源。松村英之的《Commutative Ring Theory》\cite[\S\S~8, 13--21]{Matsumura1989CommutativeRingTheory}可用于对照完备化、维数和局部结构，Hartshorne 的《Algebraic Geometry》\cite[Chapter II, \S\S~1--5]{Hartshorne1977}可供继续学习层与概形。阅读外部叙述时，请留意它采用的环、基域和有限性假设是否与眼前的命题一致；同一个定理的不同版本，条件有时正是关键所在。

\ProseParagraphBreak
我希望读者读完本卷以后，看到一组方程时，既能试着算出它的零点，也会问它的理想和商环保存了什么；看到一个局部结论时，既知道怎样在一点处验证，也愿意追问它能否成为整体结论。

\bigskip
\begin{flushright}
R. EverStarine

2026年9月24日
\end{flushright}
